\section{Duality of Feature and Sample Screening}
\label{sec:screening}

This section constructs feature screening and sample screening as
\emph{derived} operations, built from the two primitive ingredients of
Section~\ref{sec:operators}: the translation action \(T_k\) and the
coordinate restrictions \(R^F\), \(R^S\). The section has three main
results -- Decomposition, Duality, and Simultaneous Screening -- plus
one corollary (Intrinsic Complexity Reduction). Duality answers the
Duality subquestion of the paper's main question.
\noteLE{Cross-reference the actual intro label once doc1 is finalized.
Notation and dual convention as of 2026-07-03: \(T^{\mathrm
d},T^{\mathrm v},T^{\mathrm s},R^F,R^S,S^F,S^S\); \(p^*=(g^*,f^*,-A^\top,-a)\);
\(k^*=(-d,-c,-b)\); sample screening's constant \(-f_{\bar J}^*(s)\).
Section~\ref{sec:operators} (doc2) now matches.}

Throughout, \(p=(f,g,A,a)\in\mathcal FR^{\mathrm{sep}}_{m,n}\) denotes a
separable Fenchel--Rockafellar representation as in
Section~\ref{sec:operators}, with problem size \((m,n)\).

\subsection{Feature and Sample Screening as Operators}

Recall from Section~\ref{sec:operators} that an FR representation
\(p=(f,g,A,a)\) encodes the optimization problem
\[
p(x)=f(Ax)+g(x)+a.
\]
Screening simplifies this representation by eliminating a block of
variables or a block of loss terms. There are two natural ways to do
this, one on the primal side and one on the dual side, and both are
defined below from an explicit assumption on the optimal solution, with
no reference yet to translation or restriction. Only afterward, in
Theorem~\ref{thm:decomposition}, do we discover that both constructions
are the same two-step recipe: translate, then restrict.

In the literature, feature screening typically eliminates a block of
primal variables certified to vanish at the optimum. We state the
assumption slightly more generally, allowing the eliminated block to
sit at an arbitrary fixed point rather than requiring it to be zero.
Let \(I\sqcup\bar I=\{1,\ldots,n\}\) and \(z\in\mathbb R^{\bar I}\), and
suppose every optimal solution of \(p\) satisfies
\[
x_{\bar I}=z,
\]
i.e., the eliminated block of primal variables sits at the fixed point
\(z\). Writing \(x=(x_I,x_{\bar I})\), \(A=(A_I,A_{\bar I})\), and
\(g(x)=g_I(x_I)+g_{\bar I}(x_{\bar I})\) by separability, substituting
\(x_{\bar I}=z\) into the objective gives
\[
p(x_I,z)
=
f(A_Ix_I+A_{\bar I}z)+g_I(x_I)+a+g_{\bar I}(z).
\]

\begin{definition}[Feature screening]
The \emph{feature screening} of \(p\) at \((I,z)\) is the FR
representation
\[
S^{F}_{I,z}(p)
:=
\bigl(
T^{\mathrm d}_{A_{\bar I}z}f,\;
g_I,\;
A_I,\;
a+g_{\bar I}(z)
\bigr),
\]
so that \(S^F_{I,z}(p)(x_I)=p(x_I,z)\): the screened objective is
exactly the original objective with \(x_{\bar I}\) fixed at \(z\).
\end{definition}

Geometrically, feature screening collapses the eliminated block of
primal variables to the point \(z\).

\begin{example}[Feature screening for the Lasso]
\label{ex:lasso-feature}
Let \((a_i,b_i)_{i=1}^m\) be a training set with \(a_i\in\mathbb R^n\),
\(b_i\in\mathbb R\), and let \(\lambda\in\mathbb R\), \(\lambda>0\).
Writing \(A=(a_1,\ldots,a_m)^\top\in\mathbb R^{m\times n}\) and
\(b=(b_1,\ldots,b_m)^\top\), the Lasso problem
\[
\min_{x\in\mathbb R^n}
\tfrac12\|Ax-b\|_2^2+\lambda\|x\|_1
\]
is an FR representation \(p=(f,g,A,a)\) with
\[
A=A,
\qquad
f(y)=\tfrac12\|y-b\|_2^2,
\qquad
g(x)=\lambda\|x\|_1,
\qquad
a=0,
\]
separable in both \(f\) and \(g\), as required by
Section~\ref{sec:operators}.

Feature screening eliminates a block \(I\sqcup\bar I=\{1,\ldots,n\}\) of
coordinates by hypothesizing that they vanish at the optimum, the
natural hypothesis for an \(\ell_1\)-regularized problem, where
screening rules certify that some coordinates of \(x\) are exactly
zero. Take \(z=0\in\mathbb R^{\bar I}\). The two quantities entering the
definition of \(S^F_{I,z}\) both vanish,
\[
A_{\bar I}z=A_{\bar I}\cdot 0=0,
\qquad
g_{\bar I}(z)=\lambda\sum_{i\in\bar I}|0|=0,
\]
so the translation in the definition is trivial,
\(T^{\mathrm d}_{A_{\bar I}z}f=T^{\mathrm d}_0f=f\), and
\[
S^F_{I,0}(p)
=
\bigl(f,\;g_I,\;A_I,\;0\bigr),
\]
i.e., the screened representation is exactly the Lasso problem restricted
to the surviving features,
\[
\min_{x_I\in\mathbb R^{|I|}}
\tfrac12\|A_Ix_I-b\|_2^2+\lambda\|x_I\|_1,
\]
where \(A_I\in\mathbb R^{m\times|I|}\) keeps only the columns of \(A\)
indexed by \(I\).
\end{example}

We next define sample screening, the dual-side counterpart. Let
\(J\sqcup\bar J=\{1,\ldots,m\}\) and \(s\in\mathbb R^{\bar J}\), and
suppose every optimal solution of \(p\) satisfies
\[
s\in\partial f_{\bar J}(A_{\bar J}x),
\]
i.e., \(s\) is a subgradient of the eliminated loss block at the value
it takes on at optimality. By the Fenchel--Young equality, which holds
in \(f(y)+f^*(s)\geq\langle s,y\rangle\) exactly when \(s\in\partial
f(y)\),
\[
f_{\bar J}(A_{\bar J}x)
=
\langle s,A_{\bar J}x\rangle-f_{\bar J}^{*}(s).
\]
Substituting this into the objective gives
\[
p(x)
=
f_J(A_Jx)+g(x)+\langle A_{\bar J}^{\top}s,x\rangle+a-f_{\bar J}^{*}(s).
\]

\begin{definition}[Sample screening]
The \emph{sample screening} of \(p\) at \((J,s)\) is the FR
representation
\[
S^{S}_{J,s}(p)
:=
\bigl(
f_J,\;
T^{\mathrm s}_{A_{\bar J}^{\top}s}g,\;
A_J,\;
a-f_{\bar J}^{*}(s)
\bigr),
\]
so that \(S^S_{J,s}(p)(x)\) equals the original objective with
\(f_{\bar J}(A_{\bar J}x)\) replaced by its Fenchel--Young-tight value
\(\langle s,A_{\bar J}x\rangle-f_{\bar J}^{*}(s)\).
\end{definition}

Geometrically, sample screening replaces the eliminated loss block by
its exact affine value at the certified dual point \(s\).

\begin{example}[Sample screening for soft-margin SVM]
\label{ex:svm-sample}
Let \((a_i,b_i)_{i=1}^m\) be a training set with \(a_i\in\mathbb R^n\)
and \(b_i\in\{-1,+1\}\), and let \(C>0\). The soft-margin SVM problem
\[
\min_{x\in\mathbb R^n}
\tfrac12\|x\|_2^2+C\sum_{i=1}^m\max\bigl(0,\,1-b_i\langle a_i,x\rangle\bigr)
\]
seeks a hyperplane through the origin separating the two classes with as
wide a margin as possible, paying a hinge penalty for each sample that
falls inside the margin or is misclassified. Writing
\(A=(a_1,\ldots,a_m)^\top\in\mathbb R^{m\times n}\), this is an FR
representation \(p=(f,g,A,a)\) with
\[
A=A,
\qquad
f(y)=C\sum_{i=1}^m\max(0,1-b_iy_i),
\qquad
g(x)=\tfrac12\|x\|_2^2,
\qquad
a=0,
\]
separable in \(f\): the hinge loss splits sample by sample, matching the
convention that rows of \(A\) are samples.

A sample \(i\) with \(b_i\langle a_i,x\rangle>1\) lies strictly beyond
the margin on the correct side: \(f_i\) is locally zero there, so the
sample contributes nothing to the loss and cannot affect which \(x\) is
optimal. Sample screening is the operation of certifying a set of such
samples and removing them, exactly the unnecessary samples a screening
rule is meant to identify.

Let \(J\sqcup\bar J=\{1,\ldots,m\}\), and suppose every sample in
\(\bar J\) is certified safely beyond the margin, so that \(0\in\partial
f_{\bar J}(A_{\bar J}x)\) at the optimum: \(s=0\) is a valid choice in
the sample screening assumption above, since \(f_i\) is differentiable
with derivative \(0\) wherever \(b_i\langle a_i,x\rangle>1\). As in the
Lasso example, both quantities entering the definition of \(S^S_{J,s}\)
vanish at \(s=0\),
\[
A_{\bar J}^\top s=A_{\bar J}^\top\cdot 0=0,
\qquad
f_{\bar J}^*(s)=f_{\bar J}^*(0)
=\sum_{i\in\bar J}\sup_{y_i}\bigl[-C\max(0,1-b_iy_i)\bigr]=0,
\]
the last equality because \(C\max(0,1-b_iy_i)\ge0\) with equality
attained. Consequently the translation in the definition is trivial,
\(T^{\mathrm s}_{A_{\bar J}^\top s}g=T^{\mathrm s}_0g=g\), and
\[
S^S_{J,0}(p)
=
\bigl(f_J,\;g,\;A_J,\;0\bigr),
\]
i.e., the screened representation is exactly the soft-margin SVM
restricted to the surviving samples,
\[
\min_{x\in\mathbb R^n}
\tfrac12\|x\|_2^2+C\sum_{i\in J}\max\bigl(0,1-b_i\langle a_i,x\rangle\bigr),
\]
where \(A_J\in\mathbb R^{|J|\times n}\) keeps only the rows of \(A\)
indexed by \(J\). This has problem size \((|J|,n)\), down from
\((m,n)\).
\end{example}

Both definitions already display the matrix slot explicitly, so the
effect of screening on problem size needs no further argument.

\begin{corollary}[Screening reduces problem size]
\label{cor:complexity}
For every \(p=(f,g,A,a)\in\mathcal FR^{\mathrm{sep}}_{m,n}\), every
\(I\sqcup\bar I=\{1,\ldots,n\}\), and every \(J\sqcup\bar J=\{1,\ldots,m\}\),
\[
\operatorname{size}\bigl(S^F_{I,z}(p)\bigr)=(m,|I|),
\qquad
\operatorname{size}\bigl(S^S_{J,s}(p)\bigr)=(|J|,n).
\]
\end{corollary}

\begin{proof}
Immediate from the definitions above: \(S^F_{I,z}(p)\) has matrix slot
\(A_I\in\mathbb R^{m\times|I|}\), and \(S^S_{J,s}(p)\) has matrix slot
\(A_J\in\mathbb R^{|J|\times n}\).
\end{proof}

Reading intrinsic complexity as problem size, screening rules reduce
the intrinsic complexity. Theorem~\ref{thm:decomposition} below
explains why this must always happen, since translation never changes
the shape of \(A\) and restriction is the only step that can, but the
size claim itself needs nothing beyond the definitions just given.

Both definitions came from an assumption about the optimal solution, a
primal value for feature screening and a dual subgradient for sample
screening, with no reference to translation or restriction. Yet both
formulas visibly split into "shift \(f\) or \(g\), then drop
coordinates." Theorem~\ref{thm:decomposition} makes this precise.

Section~\ref{sec:operators} introduced exactly two primitives:
translation (\(T^{\mathrm d}\), \(T^{\mathrm v}\), \(T^{\mathrm s}\),
combined into \(T_k\)) and restriction (\(R^F\), \(R^S\)). The next
theorem shows that \(S^F_{I,z}\) and \(S^S_{J,s}\), defined above purely
from optimization assumptions, are always built from exactly these two
primitives and nothing else.

\subsection{Decomposition of Screening}

Look again at the two formulas above. In the feature case,
\(a+g_{\bar I}(z)\) appears together with a domain-shifted \(f\) and a
restricted \(g\); in the sample case, \(a-f_{\bar J}^*(s)\) appears
together with a slope-shifted \(g\) and a restricted \(f\). Both look
exactly like a translation of \(p\) followed by a restriction of the
result. That is the content of this theorem: the two screening
operators just defined are not primitive, they decompose.

Define the translation parameters
\[
k^{F}_{I,z}:=(A_{\bar I}z,\;g_{\bar I}(z),\;0),
\qquad
k^{S}_{J,s}:=(0,\;-f_{\bar J}^{*}(s),\;A_{\bar J}^{\top}s),
\]
and write \(T^{F}_{I,z}:=T_{k^F_{I,z}}\), \(T^{S}_{J,s}:=T_{k^S_{J,s}}\)
for the corresponding translation actions (Section~\ref{sec:operators}).
\(T^F_{I,z}\) uses a domain translation and a value translation, while
\(T^S_{J,s}\) uses a slope translation and a value translation: neither
uses all three primitives, and neither uses the same pair as the other.

\begin{theorem}[Decomposition of screening]
\label{thm:decomposition}
For every \(p=(f,g,A,a)\), every \(I\sqcup\bar I=\{1,\ldots,n\}\),
\(z\in\mathbb R^{\bar I}\), \(J\sqcup\bar J=\{1,\ldots,m\}\), and
\(s\in\mathbb R^{\bar J}\),
\[
S^{F}_{I,z}
=
R^{F}_{I}\circ T^{F}_{I,z},
\qquad
S^{S}_{J,s}
=
R^{S}_{J}\circ T^{S}_{J,s}.
\]
\end{theorem}

\begin{proof}
Since \(k^F_{I,z}\) has zero slope component,
\[
T^F_{I,z}(p)
=
\bigl(T^{\mathrm d}_{A_{\bar I}z}f,\;g,\;A,\;a+g_{\bar I}(z)\bigr).
\]
Applying \(R^F_I\), which restricts the second and third slots to
\(I\) and leaves the first and fourth untouched,
\[
R^F_I\bigl(T^F_{I,z}(p)\bigr)
=
\bigl(T^{\mathrm d}_{A_{\bar I}z}f,\;g_I,\;A_I,\;a+g_{\bar I}(z)\bigr),
\]
which is exactly \(S^F_{I,z}(p)\) from the definition above.

The sample case is the mirror computation. Since \(k^S_{J,s}\) has zero
domain component,
\[
T^S_{J,s}(p)
=
\bigl(f,\;T^{\mathrm s}_{A_{\bar J}^{\top}s}g,\;A,\;a-f_{\bar J}^{*}(s)\bigr).
\]
Applying \(R^S_J\), which restricts the first and third slots to \(J\)
and leaves the second and fourth untouched,
\[
R^S_J\bigl(T^S_{J,s}(p)\bigr)
=
\bigl(f_J,\;T^{\mathrm s}_{A_{\bar J}^{\top}s}g,\;A_J,\;a-f_{\bar J}^{*}(s)\bigr),
\]
which is exactly \(S^S_{J,s}(p)\).
\end{proof}

This is the paper's central structural fact, screening equals
translation followed by restriction, now proved rather than assumed.
Feature screening and sample screening were defined from two
unrelated-looking optimization assumptions, a primal fixed point and a
dual subgradient, and Theorem~\ref{thm:decomposition} shows they are
the same two-step construction, one primitive translation and one
primitive restriction, applied to the two different primitives of
Section~\ref{sec:operators}. This is what lets the rest of this
section's results, Duality and Simultaneous Screening, follow from
Section~\ref{sec:operators}'s Fenchel Equivariance Theorem alone,
rather than needing a separate argument for each screening type.

Theorem~\ref{thm:decomposition} also gives a second proof of
Corollary~\ref{cor:complexity}, stated already above directly from the
definitions: since translation never changes the shape of \(A\)
(Section~\ref{sec:operators}), the decomposition shows the size change
comes entirely from the restriction half, \(R^F_I\) or \(R^S_J\),
consistent with, and not needed beyond, the direct argument already
given.

\subsection{Duality of Screening}

Feature screening moves the eliminated block of \(x\) into a domain
translation of \(f\), then discards it from \(g\). Fenchel--Rockafellar
duality exchanges domain translations with slope translations and
exchanges the roles of \(f\) and \(g\) (Section~\ref{sec:operators}).
Sample screening moves an eliminated block of \(y\) into a slope
translation of \(g\), then discards it from \(f\). Matching these two
descriptions suggests that dualizing a feature screening produces
exactly a sample screening, with no other operation involved.

\begin{theorem}[Dual of screening]
\label{thm:duality}
Let \(p=(f,g,A,a)\in\mathcal FR^{\mathrm{sep}}_{m,n}\). For every
\(I\sqcup\bar I=\{1,\ldots,n\}\) and \(z\in\mathbb R^{\bar I}\),
\[
\left(S^{F}_{I,z}(p)\right)^*
=
S^{S}_{I,z}(p^*).
\]
Symmetrically, for every \(J\sqcup\bar J=\{1,\ldots,m\}\) and
\(s\in\mathbb R^{\bar J}\),
\[
\left(S^{S}_{J,s}(p)\right)^*
=
S^{F}_{J,s}(p^*).
\]
\end{theorem}

\begin{proof}
We prove the first identity; the second is the mirror computation with
the roles of \(f\) and \(g\), \(I\) and \(J\), domain and slope
translation, exchanged. The proof combines three facts from
Section~\ref{sec:operators}: Decomposition writes \(S^F_{I,z}\) as a
restriction of a translation; Restriction Equivariance dualizes the
restriction step; Fenchel Equivariance dualizes the translation step.
What remains is to check that the resulting dual translation parameter
is exactly the one that defines sample screening on \(p^*\).

By Theorem~\ref{thm:decomposition}, \(S^F_{I,z}(p)=R^F_I(T^F_{I,z}(p))\),
where \(T^F_{I,z}=T_{k^F_{I,z}}\) and
\(k^F_{I,z}=(A_{\bar I}z,\,g_{\bar I}(z),\,0)\). By Restriction
Equivariance, applied to \(q:=T^F_{I,z}(p)\),
\[
\left(S^F_{I,z}(p)\right)^*
=
\left(R^F_I(q)\right)^*
=
R^S_I(q^*)
=
R^S_I\bigl((T^F_{I,z}(p))^*\bigr).
\]
By Fenchel Equivariance,
\[
(T^F_{I,z}(p))^*
=
\bigl(T_{k^F_{I,z}}(p)\bigr)^*
=
T_{(k^F_{I,z})^*}(p^*),
\qquad
(k^F_{I,z})^*=(0,\,-g_{\bar I}(z),\,-A_{\bar I}z),
\]
using \(k^*=(-d,-c,-b)\) for \(k=(b,c,d)\).

It remains to check that \((k^F_{I,z})^*\) is exactly \(k^S_{I,z}\), the
parameter defining sample screening on \(p^*\). Writing
\(p^*=(g^*,f^*,-A^\top,-a)\), the parameter defining \(S^S_{I,z}(p^*)\)
is
\[
k^S_{I,z}
=
\bigl(0,\;-(g^*)_{\bar I}^{*}(z),\;(-A^\top)_{\bar I}^\top z\bigr).
\]
By the Lemma "Conjugate of a separable restriction,"
\((g^*)_{\bar I}=(g_{\bar I})^*\), so \((g^*)_{\bar
I}^{*}(z)=(g_{\bar I})^{**}(z)=g_{\bar I}(z)\) by biconjugation, using
that \(g_{\bar I}\) is closed, proper, convex (Lemma "Separable sums
preserve regularity"). Rows \(\bar I\) of \(-A^\top\) equal
\(-(A_{\bar I})^\top\), so \((-A^\top)_{\bar I}^\top z=-A_{\bar I}z\).
Hence
\[
k^S_{I,z}=(0,\,-g_{\bar I}(z),\,-A_{\bar I}z)=(k^F_{I,z})^*.
\]

Substituting,
\[
\left(S^F_{I,z}(p)\right)^*
=
R^S_I\bigl(T_{k^S_{I,z}}(p^*)\bigr)
=
R^S_I\bigl(T^S_{I,z}(p^*)\bigr)
=
S^S_{I,z}(p^*),
\]
the last equality by Theorem~\ref{thm:decomposition} applied to \(p^*\).
This proves the first identity.

The second identity is the mirror computation: by Decomposition,
\(S^S_{J,s}(p)=R^S_J(T^S_{J,s}(p))\); Restriction Equivariance gives
\(\left(R^S_J(T^S_{J,s}(p))\right)^*=R^F_J\bigl((T^S_{J,s}(p))^*\bigr)\);
Fenchel Equivariance gives \((T^S_{J,s}(p))^*=T_{(k^S_{J,s})^*}(p^*)\);
and the same lemma-based check, with \(f\) and \(g\), domain and slope
translation exchanged, shows \((k^S_{J,s})^*=k^F_{J,s}\), so
\(\left(S^S_{J,s}(p)\right)^*=R^F_J\bigl(T^F_{J,s}(p^*)\bigr)=S^F_{J,s}(p^*)\).
\end{proof}

Both identities are unconditional, and syntactically identical up to
exchanging \((S^F,I,z)\leftrightarrow(S^S,J,s)\): feature screening and
sample screening are literally the same construction, read on the
primal or on the dual representation. Together, they answer the
Duality subquestion of the paper's main question: feature screening on
the primal representation is exactly sample screening on the dual
representation, and the correspondence is exact, not an approximation.
There is no Commutativity theorem in this calculus: the relationship
between \(S^F\) and \(S^S\) is duality, not commutativity, and
Corollary~\ref{cor:simultaneous} below makes precise how they interact
when composed.

This theorem uses, from Section~\ref{sec:operators}: Decomposition
supplies the translate-then-restrict form of \(S^F_{I,z}\) and
\(S^S_{J,s}\); Restriction Equivariance dualizes the restriction step;
Fenchel Equivariance dualizes the translation step; and the Lemma
"Conjugate of a separable restriction" together with the Lemma
"Separable sums preserve regularity" are what let the resulting dual
translation parameter be identified with the parameter that directly
defines sample (resp.\ feature) screening on \(p^*\), the latter lemma
by letting \(g_{\bar I}\) (resp. \(f_{\bar J}\)) inherit closed, proper,
convex regularity from the individual \(g_i\) (resp. \(f_j\)), and
hence biconjugation \(g_{\bar I}^{**}=g_{\bar I}\).

Theorem~\ref{thm:duality} says a single screening dualizes to a single
screening of the other type. What happens when a feature screening and
a sample screening are both applied to the same \(p\)? Composing the
two defining formulas of the previous subsection directly and
dualizing the result depends, in general, on how the dropped samples
and the dropped features interact through \(A\): composing two
screenings on the \emph{same} representation is not order-independent
(an earlier attempt at a Commutativity theorem ran into exactly this
gap; see paper.md's Open Problems history). But duality itself
\emph{is} order-independent, in the following precise sense: screening
on \(p\) and then dualizing lands on the same representation as
dualizing first and then screening, with feature and sample screening
exchanging roles.

\begin{corollary}[Simultaneous screening]
\label{cor:simultaneous}
For every \(p=(f,g,A,a)\in\mathcal FR^{\mathrm{sep}}_{m,n}\), every
\(I\sqcup\bar I=\{1,\ldots,n\}\), \(z\in\mathbb R^{\bar I}\),
\(J\sqcup\bar J=\{1,\ldots,m\}\), and \(s\in\mathbb R^{\bar J}\),
whenever both compositions below are defined,
\[
\left(S^{S}_{J,s}\,S^{F}_{I,z}(p)\right)^*
=
S^{F}_{J,s}\,S^{S}_{I,z}(p^*),
\]
equivalently,
\[
S^{S}_{J,s}\,S^{F}_{I,z}(p)
=
\left(S^{F}_{J,s}\,S^{S}_{I,z}(p^*)\right)^*.
\]
\end{corollary}

\begin{proof}
Write \(q:=S^F_{I,z}(p)\). Applying Theorem~\ref{thm:duality} to \(q\) and \(p\), respectively, we obtain
\[
\left(S^S_{J,s}(q)\right)^*=S^F_{J,s}(q^*),
\qquad 
q^*=\left(S^F_{I,z}(p)\right)^*=S^S_{I,z}(p^*).
\]
Substituting,
\[
\left(S^S_{J,s}\,S^F_{I,z}(p)\right)^*
=
S^F_{J,s}\bigl(S^S_{I,z}(p^*)\bigr)
=
S^F_{J,s}\,S^S_{I,z}(p^*),
\]
which is the first identity. The second follows by applying
\((\cdot)^*\) to both sides and using the involution \((p^*)^*=p\)
(Proposition~\ref{prop:involution}).
\end{proof}

This is the correct replacement for a Commutativity theorem. Feature
screening and sample screening do not need to commute on the same
representation \(p\): requiring that would identify two representations
that, in general, differ by exactly the cross-block interaction
between the dropped samples and the dropped features. Instead,
composing on \(p\) and dualizing lands on the same object as composing
on \(p^*\), with \(S^F\) and \(S^S\) exchanging roles. This is what
"simultaneous screening" means in this calculus: not that the two
operators commute, but that composing them is itself compatible with
duality, exactly as a single screening is (Theorem~\ref{thm:duality}).
This is a two-line corollary of Theorem~\ref{thm:duality}, applied
twice; no new machinery is needed beyond it.
